% Manuscript excerpt; compile within main.tex for macros and references.
\subsection{Parameter counts and computational cost}
\label{app:adaptation-compute}

We distinguish parameters used in a forward prediction from parameters updated
during adaptation. Let $P_S$, $P_C$, and $P_A$ denote the parameter counts of
the pretrained local operator, composition module, and optional local
adapters. One \lddm{} prediction uses $P_S+P_C$ distinct parameters. Sharing
the local operator across $N$ subdomains and two evaluations reuses its
weights; it does not multiply the parameter count by $2N$. Composition-only
adaptation updates $P_C$ parameters while retaining the full local operator
in the forward computation.

For the Darcy source checkpoint, $P_S=3{,}957{,}666$ and
$P_C=796{,}400$, for a total of $4{,}754{,}066$ parameters. The composition
module contains 56,160 field-encoder, 141,312 spatial-encoder, and 598,928
communication parameters. Thus, composition adaptation updates 16.75\% of
the parameters used in prediction. Our matched local-adapter control inserts
low-rank updates into the local convolutional kernels. A convolution with
$C_{\rm in}$ input channels, $C_{\rm out}$ output channels, kernel size $k$,
and rank $r$ adds $r(k^2 C_{\rm in}+C_{\rm out})$ parameters. Rank 24,
capped at rank 1 for the two scalar output layers, gives $P_A=810{,}794$,
1.81\% more trainable parameters than composition adaptation.

\begin{table}[ht]
\centering
\small
\caption{\textbf{Darcy parameter accounting.} Counts are distinct parameters,
not multiplied by the number of subdomains or local calls. The local-adapter
row counts the unmerged training representation; its updates can be merged
into the original convolutional kernels for inference.}
\label{tab:darcy-parameter-accounting}
\begin{tabular}{@{}lrrr@{}}
\toprule
Adaptation & Forward parameters & Updated local & Updated composition \\
\midrule
Composition only & 4,754,066 & 0 & 796,400 \\
Local adapters only & 5,564,860 & 810,794 & 0 \\
\bottomrule
\end{tabular}
\end{table}

To separate the two sources of work, let $F_S(N,q)=Nf_S(q)$ be the forward
work of one local-operator evaluation over $N$ subdomains with $q$ grid points
each, and let $F_C(N,q)$ include token encoding, communication, and correction
decoding. With fixed input channels and trace resolution, composition depth
$L$, and token width $d$, the present dense-attention composer has
\begin{equation}
 F_C(N,q)=O\!\left(Nq+LNd^2+LN^2d\right),\qquad
 F_{\rm pred}=2F_S(N,q)+F_C(N,q)+O(Nq).
 \label{eq:composition-forward-cost}
\end{equation}
The last term covers assembly. The centered global and diagonal executions
of the composition stack change constants without changing this order. For
the fixed-resolution local CNN, the leading convolutional work is
$f_S(q)=\Theta(qP_{S,\rm conv})$, where $P_{S,\rm conv}$ counts convolutional
kernel weights; scalar-conditioning and elementwise operations add lower-order
terms for this architecture. Other local backbones can use a different
$f_S(q)$ without changing the two-call composition structure. The composition
parameter count is independent of $N$, although its dense attention costs
$O(N^2d)$ per layer.

Forward work alone does not describe adaptation. Write $B_S^x$ for the
backward work needed to differentiate a local response with respect to its
corrected input, and $B_C^\phi$ for backpropagation that computes composition
parameter gradients. For composition-only adaptation, the training-step work,
excluding common loss evaluation and optimizer overhead, is
\begin{equation}
 T_{\rm comp}=F_S^{(1)}+F_C+F_S^{(2)}
              +B_S^{x,(2)}+B_C^\phi.
 \label{eq:composition-training-cost}
\end{equation}
The first local response requires no backward pass. The second uses frozen
weights but must transmit gradients from the field loss to the predicted
input correction. Freezing therefore removes local weight-gradient
computation, not all differentiation through the local operator.

For local-adapter adaptation, let $\widetilde F_S$ include the adapter forward
work, $B_{S,A}^{(1)}$ include adapter-gradient computation through the first
local response, $B_{S,A}^{x,(2)}$ include both adapter and corrected-input
gradients through the second response, and $B_C^z$ differentiate the frozen
composer with respect to its preliminary local fields. Then
\begin{equation}
 T_{\rm local}=\widetilde F_S^{(1)}+F_C+\widetilde F_S^{(2)}
                +B_{S,A}^{(1)}+B_C^z+B_{S,A}^{x,(2)}.
 \label{eq:local-adapter-training-cost}
\end{equation}
Both local evaluations share the same adapters. For our implementation,
which forms a dense updated convolutional kernel, adapter training also
requires differentiating that effective kernel. Equal trainable parameter
counts therefore do not imply equal backward work. These expressions omit
activation-checkpoint recomputation, which trades additional forward work
for lower activation storage and is recorded separately in the ablation
protocol.

Gradient and Adam-moment storage scales as $O(P_C)$ for composition
adaptation and $O(P_A)$ for local adapters, compared with $O(P_S)$ when
fine-tuning the full local operator. The two matched arms therefore have
similar optimizer-state requirements. Total memory also includes the model
weights and activations: composition adaptation retains or recomputes the
second local evaluation for input gradients, while local adaptation must
support gradients through both local evaluations. Peak memory depends on
tile batching and checkpointing as well as parameter counts.

This separation motivates a larger pretrained local physics model with a
compact composition module. If $P_S$ grows while $P_C$ remains fixed, the
updated fraction $P_C/(P_S+P_C)$ decreases and composition optimizer storage
remains unchanged. Adaptation still incurs forward and input-gradient work
through the large local operator; it is not $O(P_C)$ computation. The present
CNN gives a moderate $P_S/P_C\approx4.97$ ratio. Scaling this advantage to
larger local backbones is a design motivation, rather than an empirical
scaling result established by these experiments.

